\chapter{Нейронные сети, которые умеют учиться: добавляем \nt{ДОФА}}
\chaptermark{Сети, которые учатся}
\label{sec:dofa}

\section{Правила игры}

Во всех схемах предыдущих глав веса ставил инженер. В организме инженера нет. Веса должны устанавливаться сами, по ходу работы, и так, чтобы сеть делала что-то полезное. Это и называется обучением.

К прежним правилам добавляется ещё одно. Синапс меняет свой вес сам, и знать он может только то, что происходит \emph{рядом с ним}: активность отправителя $x$, активность получателя $y$ и те вещества, которые до него доходят. Никто не может прислать синапсу персональную поправку.

Это исключает главный метод искусственных нейронных сетей --- обратное распространение ошибки. Там ошибка выхода проходит сеть в обратную сторону по тем же весам, в отдельную фазу после прямого прохода, и каждый вес получает свою поправку. Для этого нужны обратные провода, точно повторяющие прямые, и общий такт. Ни того ни другого в мозге не найдено, по крайней мере в той форме, что в искусственных сетях~\cite{Lillicrap2020,WhittingtonBogacz2019}.

Изменение веса будем записывать как медленный непрерывный процесс:
\begin{equation}
\frac{\dd w}{\dd t} \;=\; \eta\,\Phi(\text{то, что известно синапсу}),
\label{eq:plasticity}
\end{equation}
где $\eta$ --- скорость обучения, настолько малая, что за время одного импульса вес почти не меняется. Вся глава --- о том, какой может быть функция $\Phi$.

\section{Где хранится вес}
\label{sec:weightstore}

Что такое синапс, который «меняет свой вес»? У связи две стороны: окончание провода отправителя (\emph{пресинаптическая} сторона) и участок мембраны получателя с рецепторами (\emph{постсинаптическая} сторона), а между ними --- узкая щель. У \nt{ГЛУТ}-связей мембрана получателя обычно сидит на \emph{шипике} --- маленьком выступе входной ветви получателя. Вес связи удобно разложить на три множителя~\cite{delCastillo1954}:
\begin{empheq}[box=\keybox]{equation}
w \;=\; N_s\,p\,A_1 .
\label{eq:quantal}
\end{empheq}
Здесь $N_s$ --- число мест выброса между двумя нейронами, $p$ --- вероятность, что импульс выбросит порцию медиатора в одном месте (та самая $p$ из главы~\ref{sec:code}), $A_1$ --- ответ получателя на одну порцию, то есть сколько рецепторов её ловит. Множитель $p$ живёт у отправителя, $A_1$ --- у получателя, а $N_s$ требует обеих сторон: новые места выброса вырастают, старые исчезают, и это продолжается всю жизнь.

\emph{Решает получатель.} У типичного \nt{ГЛУТ}-синапса один из рецепторов глутамата проводит ток, только если совпали два условия: глутамат пришёл от отправителя, и мембрана получателя уже возбуждена~\cite{Nowak1984}. Это правило Хебба, встроенное в одну молекулу, --- детектор совпадений из главы~\ref{sec:glutcomputer}, поставленный на вход самого синапса. Сработав, рецептор впускает в получателя кальций, и кальций запускает изменение веса. Получатель --- единственное место, где видны сразу и вход, и выход, поэтому решение принимается у него.

\emph{Исполнить решение может любая сторона.} Самое изученное долговременное усиление идёт у получателя: в мембрану встраиваются новые рецепторы~\cite{Malinow2002}, и шипик растёт~\cite{Matsuzaki2004} --- растёт $A_1$. Но получатель может изменить и $p$: он выбрасывает вещество, которое идёт через щель назад, к отправителю, --- обратный канал из приложения~\ref{sec:othermediators}, --- и тот начинает выбрасывать меньше медиатора~\cite{WilsonNicoll2001}. А кратковременные изменения веса, истощение и облегчение из главы~\ref{sec:code}, --- это изменения $p$, которые отправитель ведёт сам по своей недавней активности.

Для инженера регистр веса разделён между двумя микросхемами. Правило обучения исполняет получатель, а записать результат он может и у себя, и у отправителя, через обратный канал. Поэтому слово «локально» в правилах этой главы означает не одну сторону связи, а всю связь целиком.

\section{Учиться без учителя}

Самое простое, что может сделать синапс, --- усилиться, когда отправитель и получатель активны вместе:
\begin{equation}
\frac{\dd w}{\dd t} \;=\; \eta\,x\,y .
\label{eq:hebb}
\end{equation}
Это правило Хебба~\cite{Hebb1949}. Оно локально и не требует часов. Но у него та же беда, что у \nt{ГЛУТ}-сети в главе~\ref{sec:glutcomputer}: положительная обратная связь. Сильный вес делает $y$ больше, большой $y$ ещё усиливает вес, и веса растут без предела. Лекарство --- отрицательная обратная связь по весу: пусть вес ещё и убывает пропорционально $y^2$. Для нейрона с входами $\mathbf x$ и линейным выходом $y=\mathbf w\cdot\mathbf x$ получаем
\begin{empheq}[box=\keybox]{equation}
\frac{\dd\mathbf w}{\dd t} \;=\; \eta\,y\,\bigl(\mathbf x - y\,\mathbf w\bigr) .
\label{eq:oja}
\end{empheq}
Правило по-прежнему локально: синапс знает свой вход $x_i$, выход $y$ и вес $w_i$.

\begin{keyblock}
\begin{proposition}[Что находит правило Хебба]
\label{prop:oja}
Правило~\eqref{eq:oja} приводит вектор весов к единичной длине вдоль направления, в котором входы меняются сильнее всего, --- к главному собственному вектору матрицы корреляций входов $C=\langle\mathbf x\mathbf x^{\mathsf T}\rangle$~\cite{Oja1982}. Нейрон учится выдавать самую выразительную одну величину, в которую можно сжать его входы.
\end{proposition}
\end{keyblock}

Доказательство. Усредним~\eqref{eq:oja} по входам: $\dd\mathbf w/\dd t=\eta\bigl(C\mathbf w-(\mathbf w^{\mathsf T}C\mathbf w)\,\mathbf w\bigr)$. Неподвижные точки --- собственные векторы $C$ единичной длины. Если $\mathbf w$ близок к собственному вектору $\mathbf e_k$ с числом $\lambda_k$, малая добавка вдоль другого вектора $\mathbf e_j$ растёт как $e^{\eta(\lambda_j-\lambda_k)t}$. Устойчива только та точка, для которой все $\lambda_j<\lambda_k$, то есть главный вектор.

Есть и вариант правила Хебба, который смотрит на время импульсов. Если импульс отправителя пришёл за $s$ миллисекунд \emph{до} импульса получателя, вес растёт на $A_+e^{-s/\tau_+}$; если через $s$ миллисекунд \emph{после} --- уменьшается на $A_-e^{-s/\tau_-}$, с $\tau_\pm$ порядка двадцати миллисекунд. Такое правило измерено на синапсах \nt{ГЛУТ}-нейронов~\cite{BiPoo1998}. Оно усиливает связи, которые \emph{могли быть причиной} ответа, и ослабляет опоздавшие (рис.~\ref{fig:stdp}). Прогоните через сеть много раз одну последовательность возбуждений, и в ней сами вырастут связи от каждого звена к следующему --- цепочка из главы~\ref{sec:glutcomputer}, собранная без инженера.

\begin{figure}[t]
\centering
\begin{tikzpicture}[>=Stealth,x=0.07cm,y=1.3cm]
\draw[->,gray!70!black] (-66,0) -- (68,0) node[right,black] {\small $s$};
\draw[->,gray!70!black] (0,-1.2) -- (0,1.35) node[above,black] {\small изменение веса};
\draw[very thick,blue!60!black] plot[domain=0.5:60,samples=50] (\x,{exp(-\x/20)});
\draw[very thick,red!70!black] plot[domain=-60:-0.5,samples=50] (\x,{-0.6*exp(\x/20)});
\draw[blue!60!black,thick] (0.5,0) -- (0.5,1);
\draw[red!70!black,thick] (-0.5,0) -- (-0.5,-0.6);
\foreach \x in {-40,-20,20,40}{\draw[gray!60!black] (\x,-0.04) -- (\x,0.04);}
\node[below,font=\scriptsize] at (20,-0.04) {$20\ms$};
\node[above,font=\scriptsize] at (-20,0.04) {$-20\ms$};
\node[align=left,font=\small,blue!60!black,anchor=west] at (22,0.95) {отправитель раньше:\\ связь усиливается};
\node[align=right,font=\small,red!70!black,anchor=east] at (-12,-0.95) {отправитель позже:\\ связь ослабевает};
\end{tikzpicture}
\caption{Правило Хебба по времени. По горизонтали --- $s=t_{\text{получ}}-t_{\text{отпр}}$, на сколько импульс получателя отстал от импульса отправителя; $\tau_\pm=20\ms$, $A_-=0.6A_+$. Окно шириной в несколько десятков миллисекунд --- того же порядка, что и окно совпадения~\eqref{eq:window}.}
\label{fig:stdp}
\end{figure}

Но все правила этого раздела учат тому, что \emph{часто}, а не тому, что \emph{полезно}. Синапс, которому известны только $x$ и $y$, не знает, принесло ли действие сок или боль. Чтобы учиться на результате, синапсу нужен третий сигнал.

\section{Третий множитель}

Третий сигнал приносит \nt{ДОФА}: дофаминовые нейроны рассылают широковещательно одно число --- ошибку предсказания награды $\delta$~\eqref{eq:rpe} (глава~\ref{sec:neurontypes}). Пусть изменение веса пропорционально произведению трёх множителей: активности отправителя, активности получателя и $\delta$.

Трудность в том, что награда приходит не тогда, когда сеть выбирала действие, а через сотни миллисекунд или секунды. К приходу $\delta$ произведение $xy$ давно равно нулю, и синапсу нужна память о том, что он участвовал. Самая простая память --- знакомая нам RC-цепь:
\begin{empheq}[box=\keybox]{equation}
\tau_e\,\frac{\dd e}{\dd t} \;=\; -e \;+\; x\,y ,
\qquad
\frac{\dd w}{\dd t} \;=\; \eta\,\delta(t)\,e(t) .
\label{eq:threefactor}
\end{empheq}
Величину $e$ называют \emph{следом}~\cite{Fremaux2016}. Совместная активность оставляет в синапсе метку, которая тает за время $\tau_e$. Если за это время приходит дофамин, вес меняется со знаком $\delta$; если нет, метка исчезает бесследно (рис.~\ref{fig:trace}). На \nt{ГЛУТ}-синапсах измерено: дофамин, пришедший в течение секунды-двух после совместной активности, превращает её в усиление связи, а пришедший позже --- нет~\cite{Yagishita2014}. Окна пластичности длиной в секунды нашли и там, где третьим сигналом служит не дофамин, а сильное возбуждение самого получателя~\cite{Bittner2017}.

\begin{figure}[t]
\centering
\begin{tikzpicture}[>=Stealth,x=7cm,y=1cm]
\fill[orange!12] (0.7,-0.2) rectangle (0.8,5.2);
% rows: x*y, delta, traces, weights
\foreach \y/\lab in {4.4/{$x\,y$},3.1/{$\delta$},1.4/{след $e$},0/{вес $w$}}{
  \draw[gray!70!black] (0,\y) -- (1.42,\y);
  \node[left,font=\small] at (-0.02,\y+0.3) {\lab};}
\draw[very thick,blue!60!black] (0,4.4) -- (0,4.9) -- (0.2,4.9) -- (0.2,4.4);
\node[right,font=\scriptsize,blue!60!black] at (0.21,4.8) {отправитель и получатель активны вместе};
\draw[very thick,orange!85!black] (0,3.1) -- (0.7,3.1) -- (0.7,3.7) -- (0.8,3.7) -- (0.8,3.1) -- (1.4,3.1);
\node[right,font=\scriptsize,orange!85!black] at (0.81,3.6) {пришёл дофамин};
% traces, each in units of its own maximum
\draw[very thick,blue!60!black] plot[domain=0:0.2,samples=20] (\x,{1.4+1.1*(1-exp(-\x))/(1-exp(-0.2))})
  plot[domain=0.2:1.4,samples=60] (\x,{1.4+1.1*exp(-(\x-0.2))});
\draw[thick,densely dashed,gray!50!black] plot[domain=0:0.2,samples=30] (\x,{1.4+1.1*(1-exp(-\x/0.05))/(1-exp(-4))})
  plot[domain=0.2:1.4,samples=80] (\x,{1.4+1.1*exp(-(\x-0.2)/0.05)});
\node[right,font=\scriptsize,blue!60!black] at (1.0,2.15) {$\tau_e=1\,\text{с}$};
\node[right,font=\scriptsize,gray!40!black] at (0.3,1.65) {$\tau_e=50\ms$};
% weights
\draw[very thick,blue!60!black] (0,0.25) -- (0.7,0.25) -- (0.8,0.8) -- (1.4,0.8) node[right,font=\scriptsize] {$\tau_e=1\,\text{с}$: вес вырос};
\draw[thick,densely dashed,gray!50!black] (0,0.2) -- (1.4,0.2) node[right,font=\scriptsize,gray!40!black] {$\tau_e=50\ms$: без изменений};
% time axis marks
\foreach \x/\l in {0/0,0.2/0.2,0.7/0.7,1.4/1.4}{\draw[gray!60!black] (\x,-0.05) -- (\x,0.05) node[below=3pt,font=\scriptsize] {\l\,с};}
\draw[<->,thick,gray!50!black] (0.2,5.35) -- node[above,font=\scriptsize,black] {задержка награды $0.5\,\text{с}$} (0.7,5.35);
\end{tikzpicture}
\caption{След по правилу~\eqref{eq:threefactor}. Совместная активность длится $0.2\,\text{с}$, дофамин приходит через полсекунды после её конца. Следы показаны в долях своего максимума. Медленный след ($\tau_e=1\,\text{с}$) к этому моменту ещё жив, быстрый ($\tau_e=50\ms$) давно растаял.}
\label{fig:trace}
\end{figure}

\begin{keyblock}
\begin{proposition}[Обучение широковещательным сигналом]
\label{prop:threefactor}
Правило~\eqref{eq:threefactor} меняет только те синапсы, которые участвовали в работе сети за последние $\tau_e$, в сторону, заданную общим сигналом $\delta$. Широковещательный сигнал вместе с локальным следом даёт адресную поправку. Задержка между действием и результатом допустима, пока она не больше $\tau_e$.
\end{proposition}
\end{keyblock}

\section{Почему это работает: шум как поиск}
\label{sec:dofagradient}

Почему правило~\eqref{eq:threefactor} ведёт к лучшему? Ответ даёт шум из главы~\ref{sec:code}. Пусть выход нейрона $y=f(u)+\xi$, где $u=\sum_jw_jx_j$, а $\xi$ --- случайное дрожание с нулевым средним и дисперсией $\sigma^2$. Награда $R$ зависит от $y$. Возьмём в качестве следа не просто $x_jy$, а его случайную часть $x_j\xi$ и в качестве третьего множителя --- ошибку $\delta=R-\bar R$, где $\bar R$ --- ожидаемая награда. При малом шуме $R\approx R\bigl(f(u)\bigr)+R'\,\xi$, поэтому
\begin{empheq}[box=\keybox]{equation}
\bigl\langle \delta\,\xi\,x_j \bigr\rangle \;=\; \sigma^2\,R'\,x_j \;=\; \frac{\sigma^2}{f'(u)}\,\frac{\partial\langle R\rangle}{\partial w_j} .
\label{eq:perturbation}
\end{empheq}
Средний шаг идёт по градиенту ожидаемой награды~\cite{Williams1992}. Никто не вычислял производных: сеть случайно отклонилась, дофамин сообщил, стало ли лучше, и участвовавшие синапсы сдвинулись в выгодную сторону (рис.~\ref{fig:perturbation}). Шум, который в главе~\ref{sec:code} мешал передавать сообщения, здесь становится поиском.

\begin{figure}[t]
\centering
\begin{tikzpicture}[>=Stealth,x=2cm,y=3cm]
\draw[->,gray!70!black] (0,0) -- (4.2,0) node[below left,font=\small,black] {выход $y$};
\draw[->,gray!70!black] (0,0) -- (0,1.2) node[left,font=\small,black] {награда $R$};
% reward hill
\draw[very thick,blue!60!black] plot[domain=0:4,samples=60] (\x,{max(0,1-((\x-3)/2.2)^2)});
\node[font=\scriptsize,blue!60!black,anchor=south] at (3,1.02) {лучший выход};
% noise around f(u)
\draw[thick,gray!60!black] plot[domain=0.6:2.4,samples=50] (\x,{0.28*exp(-((\x-1.5)/0.3)^2/2)});
\draw[densely dashed,gray!60!black] (1.5,0) -- (1.5,0.535);
\node[below,font=\small] at (1.5,0) {$f(u)$};
\node[font=\scriptsize,gray!50!black] at (1.5,-0.2) {дрожание $\xi$};
% expected reward
\draw[densely dotted,gray!60!black] (0.5,0.535) node[left,font=\scriptsize,black] {$\bar R$} -- (2.1,0.535);
% two random deviations
\fill[green!45!black] (2.1,0.833) circle (0.035);
\draw[very thick,green!45!black] (2.1,0.535) -- (2.1,0.833);
\node[font=\scriptsize,green!45!black,anchor=west,align=left] at (2.18,0.6) {$\xi>0$: лучше,\\ $\delta>0$};
\fill[red!70!black] (0.9,0.089) circle (0.035);
\draw[very thick,red!70!black] (0.9,0.535) -- (0.9,0.089);
\node[font=\scriptsize,red!70!black,anchor=east,align=right] at (0.83,0.27) {$\xi<0$: хуже,\\ $\delta<0$};
% resulting step
\draw[->,very thick,orange!85!black] (1.55,0.4) -- (1.95,0.4) node[right,font=\scriptsize,black] {шаг};
\end{tikzpicture}
\caption{Шум как поиск~\eqref{eq:perturbation}. Выход нейрона дрожит вокруг $f(u)$. Случайное отклонение вправо (зелёное) дало больше ожидаемого, $\delta>0$; влево (красное) --- меньше, $\delta<0$. В обоих случаях произведение $\delta\,\xi$ положительно, и вес сдвигается вправо, туда, где награда растёт. В среднем шаг пропорционален наклону $R'$: на вершине горки отклонения в обе стороны одинаково плохи, и шаги гасят друг друга.}
\label{fig:perturbation}
\end{figure}

\begin{keyblock}
\begin{proposition}[Спуск по градиенту без обратных проводов]
\label{prop:gradient}
Если в сети есть случайные отклонения активности, а широковещательный сигнал равен ошибке предсказания награды и в каждой обстановке в среднем равен нулю, правило~\eqref{eq:threefactor} в среднем меняет веса по градиенту ожидаемой награды. Обратные провода и отдельная фаза обучения для этого не нужны.
\end{proposition}
\end{keyblock}

Теперь понятно, почему дофамин сообщает не о награде, а об \emph{ошибке} её предсказания. Если бы третьим множителем была сама награда $R\ge0$, среднее в~\eqref{eq:perturbation} не изменилось бы, но появились бы две беды. Во-первых, к полезной части прибавилось бы слагаемое $\bar R\,\xi x_j$: в среднем ноль, но сильный шум. Во-вторых, что хуже, настоящий синапс не умеет отделять в следе $x_jy$ случайную часть от неслучайной. При данных входах $x_jf(u)$ --- одно и то же число от попытки к попытке; если $\delta$ в этой обстановке в среднем равна нулю, эта часть следа ничего не меняет, а при $\delta\ge0$ она раз за разом усиливает всё, что сеть делала, --- правильное и неправильное. Поэтому $\bar R$ должно быть ожиданием \emph{в данной обстановке}, а не средним по всей жизни. Вычислять его --- дело критика, о котором ниже.

Мы проверили это на модели. Две группы \nt{ГЛУТ}-нейронов выбирают одно из двух действий через взаимное торможение, как на рис.~\ref{fig:wta}; веса от сигнала «начали» к группам сначала равны, а шум решает, кто победит. Действие $A$ приносит сок с вероятностью $0.8$, действие $B$ --- с вероятностью $0.2$, и сок приходит через полсекунды после выбора. При $\tau_e=1\,\text{с}$ и $\delta=R-\bar R$ доля выборов $A$ за пятьсот попыток выросла с $0.65$--$0.8$ в первой сотне до $0.96$--$1.0$ в последней, а веса остались в пределах $0.85$--$1.35$. При $\tau_e=50\ms$, меньше задержки награды, сеть не научилась ничему: доля выборов $A$ осталась около половины. При $\delta=R$ без вычитания ожидания сеть тоже стала выбирать $A$, но вес победителя за те же пятьсот попыток вырос в тысячу раз и продолжал расти; а при той же $\delta=R$ и вдесятеро большей скорости обучения сеть в одном из трёх прогонов навсегда закрепила свой первый случайный выбор --- худший.

\section{Цена одного провода}

Сигнал $\delta$ один на всех, а шум, который он оценивает, --- сумма дрожаний всех $N$ нейронов, влияющих на результат. Каждый синапс видит в $\delta$ свою полезную часть и чужой шум $N-1$ соседей. Чтобы выделить своё, приходится усреднять по многим попыткам, и время обучения растёт пропорционально $N$~\cite{Werfel2005}. Обратное распространение такой платы не требует.

\begin{keyblock}
\begin{proposition}[Цена широковещания]
\label{prop:broadcastcost}
Время обучения по одному общему сигналу растёт пропорционально числу нейронов, чей шум влияет на результат. Такое обучение годится для небольших цепей, выбирающих из немногих вариантов, но не для настройки миллиардов весов.
\end{proposition}
\end{keyblock}

Мозг, по-видимому, так и делит работу: выходы \nt{ДОФА}-нейронов идут прежде всего к областям, которые выбирают действия (глава~\ref{sec:neurontypes}), а огромная \nt{ГЛУТ}-сеть коры, где весов подавляющее большинство, учится в основном местными правилами, без внешнего учителя. Раньше их сводили к правилам Хебба~\eqref{eq:oja}. Сейчас всё больше данных за то, что у коры есть собственная цель --- \emph{предсказывать} свой следующий вход, и тогда ошибка вычисляется на месте, как разница между предсказанным и пришедшим~\cite{Keller2018}: учитель встроен в саму сеть. Окна пластичности длиной в секунды, где третьим сигналом служит сильное возбуждение самого получателя~\cite{Bittner2017}, показывают, что местный сигнал ошибки у синапсов действительно бывает. Насколько это общее правило коры, --- гипотеза.

\section{Откуда берётся ошибка: критик в непрерывном времени}

Остаётся вопрос, как сами \nt{ДОФА}-нейроны вычисляют $\delta$. В программах обучения с подкреплением ошибку предсказания считают шагами $t,t+1,\dots$ В организме шагов нет, поэтому запишем её в непрерывном времени~\cite{Doya2000}. Пусть $r(t)$ --- поток награды, а $V(t)$ --- ожидание всей будущей награды, в которой более далёкая ценится меньше:
\begin{equation}
V(t) \;=\; \int_t^{\infty} e^{-(s-t)/\tau_\gamma}\,r(s)\,\dd s .
\end{equation}
Продифференцировав, получим $\dd V/\dd t=V/\tau_\gamma-r$. Невязка этого равенства и есть ошибка предсказания:
\begin{empheq}[box=\keybox]{equation}
\delta(t) \;=\; r(t) \;+\; \frac{\dd V}{\dd t} \;-\; \frac{V}{\tau_\gamma} .
\label{eq:ctd}
\end{empheq}
Постоянная $\tau_\gamma$ --- горизонт планирования, непрерывный аналог коэффициента $\gamma$; по гипотезе, его задаёт \nt{СЕРО} (раздел~\ref{sec:horizon}), и тогда~\eqref{eq:patience} --- правило, сколько стоит ждать. Проверим три случая (рис.~\ref{fig:ctdcases}). Загорелась лампочка, обещающая сок: $V$ скачком выросло, $\dd V/\dd t$ большое, всплеск. Пришёл обещанный сок: $r>0$, но ожидание в тот же момент исчерпалось, $\dd V/\dd t\approx-r$, и всплеска нет. Сок не пришёл: ожидание исчезло без награды, $\dd V/\dd t<0$, провал.

\begin{figure}[t]
\centering
\begin{tikzpicture}[>=Stealth,x=1.15cm,y=1cm]
\foreach \col/\ttl/\lampon/\juice in {0/{неожиданный сок}/0/1, 1/{обещанный сок}/1/1, 2/{сок не пришёл}/1/0}{
 \begin{scope}[xshift=\col*4.6cm]
  \node[font=\small] at (1.5,3.55) {\ttl};
  \foreach \yy in {2.4,1.2,0}{\draw[gray!60!black] (0,\yy) -- (3.1,\yy);}
  % reward r
  \ifnum\juice=1
    \fill[green!45!black] (1.95,2.4) rectangle (2.05,2.85);
    \pic[scale=0.55] at (2.0,3.1) {juice};
  \else
    \pic[scale=0.55] at (2.0,3.1) {nojuice};
  \fi
  % expectation V and error delta
  \ifnum\lampon=1
    \pic[scale=0.55] at (1.0,3.1) {lamp};
    \draw[very thick,blue!60!black] (0,1.2) -- (1,1.2) -- (1,1.45)
      plot[domain=1:2,samples=20] (\x,{1.2+0.25*exp((\x-1)/1.5)}) -- (2,{1.2+0.25*exp(1/1.5)}) -- (2,1.2) -- (3.1,1.2);
    \draw[very thick,red!70!black] (0,0) -- (0.95,0) -- (0.95,0.5) -- (1.05,0.5) -- (1.05,0) -- (3.1,0);
    \ifnum\juice=0
      \draw[very thick,red!70!black] (1.95,0) -- (1.95,-0.45) -- (2.05,-0.45) -- (2.05,0);
      \node[font=\scriptsize,red!70!black,anchor=west] at (2.1,-0.35) {провал};
    \else
      \node[font=\scriptsize,gray!50!black,anchor=north,align=center] at (2.0,-0.08) {$r$ и спад $V$\\ гасят друг друга};
    \fi
  \else
    \draw[very thick,blue!60!black] (0,1.2) -- (3.1,1.2);
    \draw[very thick,red!70!black] (0,0) -- (1.95,0) -- (1.95,0.5) -- (2.05,0.5) -- (2.05,0) -- (3.1,0);
  \fi
  \ifnum\col=0
    \node[font=\small,anchor=east] at (-0.1,2.55) {$r$};
    \node[font=\small,anchor=east] at (-0.1,1.35) {$V$};
    \node[font=\small,anchor=east] at (-0.1,0.1) {$\delta$};
  \fi
 \end{scope}}
\end{tikzpicture}
\caption{Три случая для ошибки~\eqref{eq:ctd}; сверху вниз --- награда $r$, ожидание $V$ и ошибка $\delta$. Слева ожидания нет, и сок целиком --- неожиданность. В середине лампочка скачком поднимает $V$: это скачок $\dd V/\dd t$, всплеск $\delta$ приходится на лампочку. Дальше $V$ растёт ровно так, как требует горизонт, и $\delta=0$; в момент сока награда $r$ и спад $V$ гасят друг друга. Справа $V$ падает без награды --- провал. Это те же всплеск, перенос и провал, что на рис.~\ref{fig:rpe}, но теперь видно, откуда они берутся.}
\label{fig:ctdcases}
\end{figure}

Что нужно для~\eqref{eq:ctd} из схем, которые у нас уже есть? Три вещи.

\emph{Сама оценка $V$.} Её выдаёт группа \nt{ГЛУТ}-нейронов --- критик, --- веса которой учатся тем же правилом~\eqref{eq:threefactor} с тем же $\delta$: если $\delta>0$, ожидание было занижено, и связи, активные перед этим, усиливаются.

\emph{Производная $\dd V/\dd t$.} Её вычисляет любая схема, отвечающая на изменения, а не на уровень: прямое возбуждение вместе с задержанным торможением через \nt{ГАМК}-посредника, как в детекторе направления главы~\ref{sec:gaba}, или истощающийся синапс~\eqref{eq:depression} из главы~\ref{sec:code}.

\emph{Знак в одном проводе.} Ошибка бывает и положительной, и отрицательной, а частота импульсов --- только положительной. Решение то же, что у \nt{СЕРО}-нейрона в главе~\ref{sec:serocomputer}: \nt{ДОФА}-нейрон --- ритмоводитель. Его ровные несколько импульсов в секунду означают $\delta=0$, пачка --- $\delta>0$, пауза --- $\delta<0$. Отрицательной ошибке остаётся меньше места: частоту нельзя опустить ниже нуля. У настоящих \nt{ДОФА}-нейронов провалы действительно мельче всплесков~\cite{Bayer2005}.

Что дофамин ведёт себя как $\delta$, измерено. Как именно мозг вычисляет $\dd V/\dd t$, --- гипотеза.

\section{Актёр и критик}

Соберём всё вместе (рис.~\ref{fig:actorcritic}). Нейроны обстановки возбуждают группы действий, выбирающие победителя через \nt{ГАМК} (утверждение~\ref{prop:wta}), --- это \emph{актёр} --- и критика, выдающего $V$~\cite{SuttonBarto2018}. \nt{ДОФА}-нейрон получает награду $r$ и $V$, вычисляет~\eqref{eq:ctd} и рассылает $\delta$ всем синапсам актёра и критика.

\begin{figure}[t]
\centering
\begin{tikzpicture}[>=Stealth,
  nrn/.style={circle,draw,very thick,fill=blue!6,minimum size=0.9cm,inner sep=0pt,font=\small},
  gab/.style={circle,draw,very thick,fill=red!10,minimum size=0.6cm,inner sep=0pt},
  dofa/.style={circle,draw,very thick,double,double distance=1.2pt,fill=orange!15,minimum size=1.35cm,inner sep=0pt,font=\scriptsize},
  inh/.style={-{Bar[width=7pt]},thick,red!70!black},
  bc/.style={->,thick,densely dashed,orange!85!black}]
\node[nrn] (S1) at (0,2.4) {$x_1$};
\node[nrn] (S2) at (0,1.2) {$x_2$};
\node[nrn] (S3) at (0,0) {$x_3$};
\node[left=0.15cm,align=right,font=\small] at (-0.5,1.2) {обстановка};
\node[nrn] (A) at (4,2.6) {$A$};
\node[nrn] (B) at (4,1.0) {$B$};
\node[gab] (G) at (5.2,1.8) {};
\draw[->,thick] (A) -- (G); \draw[->,thick] (B) -- (G);
\draw[inh] (G) to[bend left=35] (A);
\draw[inh] (G) to[bend right=35] (B);
\node[above,font=\small] at (4.6,3.05) {актёр: выбор действия};
\node[nrn] (V) at (4,-1.1) {$V$};
\node[below,font=\small] at (4,-1.6) {критик};
\foreach \s in {S1,S2,S3}{\foreach \t in {A,B,V}{\draw[->,gray!60!black] (\s) -- (\t);}}
\node[dofa] (D) at (8.0,-1.1) {\nt{ДОФА}};
\draw[->,thick] (V) -- node[above,font=\scriptsize] {$V,\ \dd V/\dd t$} (D);
\draw[->,thick,green!45!black] (10.0,-1.1) node[right,black,font=\small] {награда $r$} -- (D);
\pic[scale=1.1] at (10.6,-0.5) {juice};
\draw[bc] (D.north) -- (8.0,4.0) -- node[above,font=\small,black] {$\delta$ --- всем синапсам актёра и критика} (2.2,4.0) -- (2.2,3.0);
\draw[bc] (D.south) -- (8.0,-2.4) -- (2.2,-2.4) -- (2.2,-0.6);
\end{tikzpicture}
\caption{Сеть, которая учится выбирать действия. Серые стрелки --- \nt{ГЛУТ}-синапсы с меняющимися весами; красный кружок --- \nt{ГАМК}-нейрон; двойной кружок --- \nt{ДОФА}-ритмоводитель, вычисляющий $\delta$ по~\eqref{eq:ctd}; штриховые стрелки --- широковещательная рассылка $\delta$.}
\label{fig:actorcritic}
\end{figure}

Знак действия медиатора выбирает получатель (утверждение~\ref{prop:receptor}), и в области выбора действий часть нейронов несёт дофаминовые рецепторы одного семейства, часть --- другого. В опытах на срезах дофамин вместе с совместной активностью у первых усиливает синапсы, у вторых --- ослабляет~\cite{Shen2008}; в модели это значит, что у вторых знак $\delta$ в~\eqref{eq:threefactor} перевёрнут. Первые учатся тому, что \emph{стоит делать}, вторые --- тому, чего \emph{делать не стоит}.

\section{Итог}

\begin{table}[t]
\centering
\footnotesize
\setlength{\tabcolsep}{4pt}
\renewcommand{\arraystretch}{1.15}
\begin{tabular}{>{\raggedright}p{2.6cm}>{\raggedright}p{2.7cm}>{\raggedright}p{3.1cm}>{\raggedright\arraybackslash}p{5.0cm}}
\toprule
Правило & Что знает синапс & Чему учится сеть & Что известно \\
\midrule
Хебб по частотам~\eqref{eq:oja} & $x$, $y$, свой вес & сжатию входов: главным направлениям & усиление при совместной активности измерено; механизм ограничения весов --- предмет исследований \\
Хебб по времени & порядок импульсов $x$ и $y$ в пределах $\sim20\ms$ & причинным связям, последовательностям & измерено на \nt{ГЛУТ}-синапсах \\
три множителя~\eqref{eq:threefactor} & $x$, $y$, след $e$, общий сигнал $\delta$ & действиям, которые приносят награду & влияние дофамина в течение секунд после активности измерено; как главный механизм обучения выбору --- хорошо обоснованная гипотеза \\
критик~\eqref{eq:ctd} & то же & ожиданию награды $V$ & сходство дофамина с $\delta$ измерено; схема, вычисляющая $\dd V/\dd t$, --- гипотеза \\
обратное распространение ошибки & персональную поправку для каждого веса & всему, но нужны обратные провода и такт & в мозге не найдено в этой форме; ищут приближения к нему \\
\bottomrule
\end{tabular}
\caption{Правила обучения в сети из \nt{ГЛУТ}-, \nt{ГАМК}- и \nt{ДОФА}-нейронов.}
\label{tab:learning}
\end{table}

Правила обучения сведены в табл.~\ref{tab:learning}. \nt{ГЛУТ} несёт данные, \nt{ГАМК} управляет потоком, а \nt{ДОФА} решает, какие изменения в сети оставить.

\section{Задачи}

\begin{exercise}[\lvlA]
\label{ex:06:1}
\emph{Сколько живёт след.} В опыте рис.~\ref{fig:trace} совместная активность $xy=1$ длится $0.2\,\text{с}$, начиная с $e=0$, а дофамин приходит через $0.5\,\text{с}$ после её конца.
(а)~Найдите по~\eqref{eq:threefactor} значение следа в момент прихода дофамина при $\tau_e=1\,\text{с}$ и при $\tau_e=50\ms$ и их отношение.
(б)~При какой наибольшей задержке след при $\tau_e=1\,\text{с}$ ещё не меньше $10\,\%$ своего значения в конце совместной активности?
(в)~При фиксированных длительности активности $T_a$ и задержке $d$ найдите $\tau_e$, при котором след в момент дофамина наибольший. Вычислите его для $T_a=0.2\,\text{с}$ и для $T_a=1\,\text{с}$ при $d=0.5\,\text{с}$. Почему слишком длинный след тоже плох?
\end{exercise}

\begin{exercise}[\lvlA]
\label{ex:06:2}
\emph{Дрейф правила Хебба по времени.} Отправитель и получатель выдают независимые пуассоновские потоки с частотами $\nu_x$ и $\nu_y$; каждая пара импульсов меняет вес по окну рис.~\ref{fig:stdp}.
(а)~Докажите, что средняя скорость изменения веса равна $\nu_x\nu_y\,(A_+\tau_+-A_-\tau_-)$.
(б)~Вычислите её при $\tau_\pm=20\ms$, $A_-=0.6A_+$, $\nu_x=\nu_y=10$ импульсов в секунду. На сколько $A_+$ вырастет вес за час совершенно случайной активности?
(в)~При каком $\tau_-$ (при $A_-=0.6A_+$) дрейф исчезает? Почему это условие нужно выполнять точно, если правило должно учить только причинным связям?
\end{exercise}

\begin{exercise}[\lvlB]
\label{ex:06:3}
\emph{Доказательство утверждения~\ref{prop:oja} до конца.} Усреднённое правило: $\dd\mathbf w/\dd t=\eta\bigl(C\mathbf w-(\mathbf w^{\mathsf T}C\mathbf w)\mathbf w\bigr)$.
(а)~Покажите, что $\dd\|\mathbf w\|^2/\dd t=2\eta\,(\mathbf w^{\mathsf T}C\mathbf w)(1-\|\mathbf w\|^2)$, и выведите, что длина вектора весов стремится к единице.
(б)~Линеаризуйте правило в точке $\mathbf e_k$. Покажите, что отклонение вдоль $\mathbf e_j$, $j\ne k$, растёт как $e^{\eta(\lambda_j-\lambda_k)t}$, а отклонение длины затухает со скоростью $2\eta\lambda_k$.
(в)~Для $C=\begin{pmatrix}2&1\\1&2\end{pmatrix}$ найдите предел $\mathbf w$ и самую медленную постоянную времени сходимости.
(г)~Частоты импульсов положительны, и у входов есть среднее: $\mathbf x=\mathbf m+\mathbf n$, $\langle\mathbf n\rangle=0$, $\langle\mathbf n\mathbf n^{\mathsf T}\rangle=\Sigma$. Пусть $\mathbf m=(\mu,0)$, $\Sigma=\operatorname{diag}(s_1^2,s_2^2)$. При каком условии нейрон выучит направление $\mathbf e_1$? Проверьте его при $\mu=1$, $s_1=0.1$, $s_2=0.5$ и уточните формулировку утверждения~\ref{prop:oja}.
\end{exercise}

\begin{exercise}[\lvlA]
\label{ex:06:4}
\emph{Непрерывный критик как предел шагов.}
(а)~Разбейте время на шаги $\Delta$, награду за шаг возьмите $r(t)\Delta$, а $\gamma=e^{-\Delta/\tau_\gamma}$. Покажите, что $\delta_t/\Delta$ из~\eqref{eq:rpe} при $\Delta\to0$ переходит в~\eqref{eq:ctd}.
(б)~Лампочка зажигается в непредсказуемый момент $t_c$, сок размера $R$ (короткий импульс площади $R$) приходит в момент $t_c+L$. Для выученного ожидания найдите $V(t)$ и $\delta(t)$ на всём опыте, включая случай, когда сок не пришёл. Покажите, что между лампочкой и соком $\delta\equiv0$.
(в)~Вычислите площадь всплеска на лампочку в долях $R$ при $L=1\,\text{с}$ для $\tau_\gamma=2\,\text{с}$ и $\tau_\gamma=4\,\text{с}$. Какое проверяемое предсказание отсюда следует для гипотезы раздела~\ref{sec:horizon}?
\end{exercise}

\begin{exercise}[\lvlB]
\label{ex:06:5}
\emph{Откуда берётся цена широковещания.} Пусть $N$ нейронов дрожат независимо, $\xi_i\sim\mathcal N(0,\sigma^2)$, и при малом шуме $\delta=\sum_ia_i\xi_i$, где $a_i=\partial R/\partial y_i$. Синапс $j$-го нейрона оценивает градиент величиной $\hat g_j=\delta\,\xi_jx_j$.
(а)~Найдите $\langle\hat g_j\rangle$ и $\operatorname{Var}\hat g_j$.
(б)~При одинаковых $a_i$ найдите отношение сигнала к шуму одной попытки и число попыток, нужное для заданной точности. Выведите отсюда утверждение~\ref{prop:broadcastcost}.
(в)~Пусть вместо ошибки рассылается сама награда, $\delta=\bar R+\sum_ia_i\xi_i$. Во сколько раз вырастет дисперсия? Сравните с рассуждением после утверждения~\ref{prop:gradient}.
(г)~Пусть один нейрон ($N=1$) учится задаче за $100$ попыток, а одна попытка занимает $5\,\text{с}$. Оцените время обучения при $N=10^4$ и при $N=10^{10}$.
\end{exercise}

\begin{exercise}[\lvlB]
\label{ex:06:6}
\emph{Проектирование: схема, вычисляющая $\delta$.} \nt{ДОФА}-ритмоводитель с фоновой частотой $\rho_0$ получает награду $r$, оценку $V$ прямым возбуждающим путём с весом $k_1$ и её сглаженную копию $\bar V$, $\tau_d\,\dd\bar V/\dd t=-\bar V+V$, тормозящим путём с весом $k_2$. Его частота $\rho=\rho_0+\lambda\,\hat\delta$, $\hat\delta=r+k_1V-k_2\bar V$, ограничена снизу нулём.
(а)~Выберите $k_1$, $k_2$ так, чтобы для медленных $V$ было $\hat\delta\approx$~\eqref{eq:ctd}.
(б)~Выученное ожидание растёт как $V=V_0e^{t/\tau_\gamma}$, и истинная $\delta=0$. Найдите точно, что выдаёт схема. Какое $\tau_d$ нужно, чтобы ложная ошибка не превышала $5\,\%$ от $V/\tau_\gamma$ при $\tau_\gamma=2\,\text{с}$?
(в)~$V$ скачком выросло на $\Delta V$ (лампочка). Найдите $\hat\delta(t)$ и покажите, что площадь всплеска равна $\Delta V$ при любом $\tau_d$.
(г)~Годится ли в качестве сглаживающего звена истощающийся синапс~\eqref{eq:depression} с $U=0.5$, $\tau_{\text{вос}}=0.8\,\text{с}$ при частоте $20$ импульсов в секунду? Какой будет ложная ошибка?
(д)~Пусть $\rho_0=5$, а наибольшая частота $30$ импульсов в секунду. Каков диапазон представимых $\hat\delta$ в каждую сторону? Если $\hat\delta$ гауссова с нулевым средним и $\lambda\sigma_\delta=\rho_0$, найдите среднее переданного сигнала после обрезки нулём. Чем опасно это смещение для правила~\eqref{eq:threefactor}?
\end{exercise}

\begin{exercise}[\lvlB]
\label{ex:06:7}
\emph{Моделирование: предельная задержка награды.} Работайте с копией \texttt{models/\allowbreak dofa\_learning.py}.
(а)~Добавьте в \texttt{run()} параметр задержки награды $d$ и сделайте длину попытки равной $T_{\text{cue}}+d+T_{\text{rew}}+0.4\,\text{с}$ (при $d=0.5\,\text{с}$ модель должна совпасть с исходной).
(б)~При $d=0.5\,\text{с}$, $\eta=0.05$, $500$ попытках и шести зёрнах постройте долю выборов $A$ в последней сотне попыток в зависимости от $\tau_e\in\{0.05,0.1,0.2,0.5,1,2,4\}\,\text{с}$.
(в)~При $\tau_e=1\,\text{с}$ постройте ту же долю в зависимости от $d\in\{0.25,0.5,1,2,3\}\,\text{с}$.
(г)~Вычислите для каждого прогона множитель $F=(1-e^{-T_{\text{cue}}/\tau_e})\,e^{-d/\tau_e}$ --- долю следа, доживающую до награды (задача~\ref{ex:06:1}), и постройте все точки как функцию $F$. Проверьте утверждение~\ref{prop:threefactor}: где проходит граница обучения и почему доля выборов $A$ убывает и при слишком большом $\tau_e$?
\end{exercise}

\begin{exercise}[\lvlC]
\label{ex:06:8}
\emph{Можно ли обойти цену широковещания.} Модель: $N$ нейронов, выход $y_i=w_i+\xi_i$, $\xi_i\sim\mathcal N(0,\sigma^2)$, награда $R=-\sum_i(y_i-w_i^*)^2$, ошибка $\delta=R-\langle R\rangle$ (ожидание при текущих весах), правило $\Delta w_i=\eta\,\delta\,\xi_i$ после каждой попытки. Пусть $E=\sum_i(w_i-w_i^*)^2$.
(а)~Докажите, что $\langle\Delta E\rangle=-4\eta\sigma^2E+\eta^2\bigl[4\sigma^4(N+2)E+\sigma^6N(2N+8)\bigr]$. Найдите лучшее $\eta$ и покажите, что вдали от шумового дна $E$ убывает в $e$ раз за $N+2$ попытки.
(б)~Пусть в каждой попытке дрожат только $m$ случайно выбранных нейронов из $N$. Покажите, что постоянная времени равна $N(m+2)/m$. Помогает ли разреженный поиск?
(в)~Проверьте (а) и (б) моделью: $\sigma=0.1$, начальные $w_i-w_i^*=1$, $\eta$ --- половина от $1/(\sigma^2(m+2))$; измерьте число попыток до $E<0.1E_0$ при $N=4,16,64,256$ и $m=1,4,N$ и сравните с теорией.
(г)~Открытая часть. Линейный рост с $N$ неизбежен для любого правила, у которого на входе одно число на попытку? Какие изменения условий его ломают? Рассмотрите хотя бы: несколько проводов ошибки (глава~\ref{sec:dofaalt}), местную ошибку предсказания входа, шум в весах вместо шума в нейронах. Оцените, какой из вариантов позволил бы учить $10^{10}$ весов за разумное время.
\end{exercise}
